% !TEX root = ../main.tex

\chapter{实验脚本与日志模板}

本附录在正式版中应补充关键脚本入口、命令行参数模板、日志字段样例与目录组织方式，以支撑“统一脚本、统一日志、统一口径”的研究方法。相关脚本与复核记录可与内部材料统一维护和版本化管理\citep{internal2026record}。

\section{待补充内容}

\begin{itemize}
  \item 关键训练、建索引、查询与评测脚本入口。
  \item 常用命令行参数模板与默认值说明。
  \item 日志字段示例及字段解释。
  \item 工程目录组织方式与结果归档规范。
\end{itemize}


\chapter{图表补充说明}

本附录对正文图表中反复出现的变量命名、单位和比较口径作统一补充，便于阅读时快速核对不同章节之间的含义是否一致。

\section{图表口径说明}

\begin{itemize}
  \item 图中的 QPS、平均时延和 p99 时延均指对应实验场景下的端到端统计结果；若未特别说明，时延单位统一为 ms。
  \item 单卡结果图主要用于比较数据路径优化前后的链路差异，多 GPU 结果图主要用于比较 replicated 主方案、GPU 侧归并兜底路径以及共享路径冲突条件下的扩展差异。
  \item 真实集结果用于有效性验证，1M 及以上合成数据结果用于系统性能与扩展分析，二者不直接比较绝对吞吐数值。
  \item 正文结果图中的方案名称均对应具体验证场景或系统路径，不再使用含义不明确的编号标签。
\end{itemize}

\chapter{缩写、变量与配置对照}

本附录列出正文中频繁出现的关键变量、配置项及其含义，便于后续统一实验脚本和结果复核。

\renewcommand{\arraystretch}{1.2}
\small
\begin{tabularx}{\textwidth}{>{\raggedright\arraybackslash}p{0.34\textwidth}>{\raggedright\arraybackslash}X}
变量或配置项 & 含义 \\
\texttt{batch\_size} & 单次提交到候选召回或精排路径中的批处理大小。 \\
\texttt{rerank\_budget} & 进入高精度精排阶段的候选数量预算。 \\
\texttt{nlist} & 索引聚类桶数量或粗排分区规模参数。 \\
\texttt{nprobe} & 查询阶段实际访问的粗排桶数量。 \\
\texttt{device\_id} & 当前任务绑定的 GPU 设备标识。 \\
\texttt{replica\_mode} & 多 GPU 复制式多实例执行模式开关。 \\
\texttt{bounded\_batch\_size} & 在时延约束下允许形成的微批上限。 \\
\end{tabularx}
